

%قبل از شروع مسئله ترجیحاً چیزی ننویسید. اما اگر خواستید بنویسید، در قالب لیستی (فقرات و شمارش) بنویسید.
%برای اضافه کردن نمره، بعدی از شروع‌مسئله، [x] را اضافه کنید و به جای x  نمره مورد نظر را قرار دهید. برای زیرمسئله‌ها هم به همین صورت است.
%\شروع‌مسئله[10]
\شروع‌مسئله []
%\متن‌سیاه{سوال عملی اول:}\\
%\متن‌سیاه{دراپ‌تست}\\
علیرضا یک سلاح جدید ساخته و می‌خواهد آن را تست کند. او در محله‌ای زندگی می‌کند که دارای \متن‌لاتین{n} ساختمان است که روی یک خط قرار دارند و از چپ به راست با اعداد 1 تا \متن‌لاتین{n} شماره‌گذاری شده‌اند و ساختمان \متن‌لاتین{i} ام \متن‌لاتین{h[i]} طبقه دارد. \\
طرز کار این سلاح به این شکل است که روی پشت بام یکی از ساختمان‌ها مستقر می‌شود و به سمت راست شلیک می‌کند. تیر در مسیر افقی به حرکت خود ادامه می‌دهد تا به یک ساختمان برخورد کند (که اولین ساختمانی است که ارتفاع آن از ساختمانی که سلاح روی آن قرار گرفته بیشتر مساوی است و سمت راست سلاح قرار دارد) . \\
وقتی تیر به ساختمان برخورد کرد هسته‌ی تیر شکافته شده و یک لیزر از آن خارج می‌شود و از ساختمانی که به آن برخورد کرده بالا می‌رود و روی پشت یام آن (یعنی درست روی سقف طبقه‌ی آخر) قرار می‌گیرد. سپس به سمت راست به شکل افقی لیزر پرتاب می‌کند و هر طبقه‌ از ساختمان‌های سمت راست که ارتفاع بیشتری از لیزر داشته‌باشد (یعنی طبقاتی که ارتفاع آن‌ها از ارتفاع ساختمان لیزر بیشتر است) را نابود میکند.\\
به او کمک کنید بدون خراب کاری متوجه شود به ازای هر ساختمان اگر سلاح را روی آن ساختمان قرار دهد و شلیک کند میزان خرابی ایجاد شده چند طبقه خواهد بود.\\
توجه‌: در سوالات عملی حق استفاده از کد‌های آماده‌ی موجود در اینترنت و کتاب‌خانه‌های آماده‌ی پایتون را ندارید و باید خودتان تمام موارد را پیاده سازی کنید.\\
توجه‌:  در این سوال باید از درخت AVL استفاده کنید.

\متن‌سیاه{ورودی:}\\
در خط اول ورودی عدد صحیح \متن‌لاتین{n} آمده که تعداد ساختمان‌های محله را مشخص می‌کند.
در خط بعدی \متن‌لاتین{n} عدد آمده است که عدد \متن‌لاتین{i} ام نشان دهنده‌ی تعداد طبقات ساختمان \متن‌لاتین{i} ام یعنی \متن‌لاتین{h[i]} است.\\
\متن‌سیاه{خروجی:}\\
خروجی باید شامل \متن‌لاتین{n} عدد باشد که عدد \متن‌لاتین{i} ام نشان دهنده‌ی تعداد طبقات خراب شده در حالتی است که سلاح روی ساختمان \متن‌لاتین{i} ام قرار گرفته و شلیک کند.\\
\متن‌سیاه{محدودیت‌ها:}\\
$1 \leq n \leq 10^{5}$\\
$1 \leq h[i] \leq 10^{9}$\\

\متن‌سیاه{ورودی و خروجی نمونه:}\\

تست اول
\begin{latin}
	\begin{tabularx}{\textwidth}{|X | X|}
		\hline
		\textbf{stdin} &
		\textbf{stdout}\\
		\hline
		9 & 17 12 8 8 5 5 1 0 0\\
		1 2 3 1 4 2 5 7 8 & \\
		\hline
	\end{tabularx}
\end{latin}

تست دوم
\begin{latin}
	\begin{tabularx}{\textwidth}{|X | X|}
		\hline
		\textbf{stdin} &
		\textbf{stdout}\\
		\hline
		4 & 3 1 0 0\\
		1 2 3 4 & \\
		\hline
		
	\end{tabularx}
\end{latin}
%}
% صورت اصلی مسئله را در این‌جا بنویسید.
%اگر مسئله چند قسمت داشت به این صورت عمل کنید:
%\شروع{زیرمسائل}
%\زیرمسئله[0] % مسئلۀ ۱آ
%\زیرمسئله % مسئلۀ ۱ب
%\زیرمسئله % مسئلۀ ۱پ
%\زیرمسئله % مسئلۀ ۱ت
%\پایان{زیرمسائل}



%%%%%%%%%%%%%%%%%%%%%%%%%%%%%%%%%%%%%%%%%%%%%%%%%%%%%
% برخی دستورات مفید برای تمرین‌ها. %
%%%%%%%%%%%%%%%%%%%%%%%%%%%%%%%%%%%%%%%%%%%%%%%%%%%%%

% فرمول‌های نمونه:
%$f(x) = \Theta(x)$
%$T(x, y) \leq \log(x) + 2^y + \binom{2n}{n}$
% {\tt code\_function}


% لیست موارد بدون شماره‌گذاری:
%\شروع{فقرات}
%    \فقره مورد اول 
%        \شروع{فقرات}
%            \فقره مورد اول داخلی 
%                \شروع{فقرات}
%                    \فقره اولین مورد خیلی داخلی 
%                    \فقره دومین مورد خیلی داخلی 
%                \پایان{فقرات}
%            \فقره مورد دوم داخلی 
%        \پایان{فقرات}
%    \فقره مورد دوم 
%\پایان{فقرات}

% لیست موارد شماره‌گذاری شده:
%\شروع{شمارش}
%    \فقره یک مورد 
%    \فقره یک مورد دیگر 
%    \فقره باز هم یک مورد دیگر
%\پایان{شمارش}

% برای نوشتن یک فرمول چند خطی که یک شماره بخوره. برای حالت‌های دیگر به سایت زیر مراجعه کنید:
% https://www.overleaf.com/learn/latex/Aligning_equations_with_amsmath
%\begin{align} 
%    \begin{split}
%        (x+y)^3 &= (x+y)^2(x+y) \\
%                &= (x^2+2xy+y^2)(x+y) \\
%                &= (x^3+2x^2y+xy^2) + (x^2y+2xy^2+y^3) \\
%                &= x^3+3x^2y+3xy^2+y^3
%    \end{split}                                 
%\end{align}

% نوشتن ماتریس و آرایه دوبعدی
%\begin{align}
%    A = 
%    \begin{bmatrix}
%        A_{11} & A_{21} \\
%        A_{21} & A_{22}
%    \end{bmatrix}
%\end{align}

% اضافه کردن تصویر با شرح. برای اضافه کردن تصویر بدون شرح، فقط دستور درج‌تصویر را استفاده کنید.
%\شروع{شکل}[h!]
%\تنظیم‌ازوسط
%  \درج‌تصویر[پهنا=0.5\پهنای‌متن]{img.jpg}
%\شرح{شرح تصویر}
%\برچسب{تصویر:۱}
%\پایان{شکل}
%\پایان{زیرمسائل}

%کد زیر را ببینید.
%%\lstinputlisting{source_filename.py}
%\begin{latin}
%\begin{lstlisting}
%# Import the modules
%import sys
%import random
%
%ans = True
%
%while ans:
%question = raw_input("Ask the magic 8 ball a question: (press enter to quit) ")
%
%answers = random.randint(1,8)
%
%if question == "":
%sys.exit()
%
%elif answers == 1:
%print "It is certain"
%
%elif answers == 2:
%print "Outlook good"
%
%elif answers == 3:
%print "You may rely on it"
%
%elif answers == 4:
%print "Ask again later"
%
%elif answers == 5:
%print "Concentrate and ask again"
%
%elif answers == 6:
%print "Reply hazy, try again"
%
%elif answers == 7:
%print "My reply is no"
%
%elif answers == 8:
%print "My sources say no"
%\end{lstlisting}
%\end{latin}

